% ============================================================
%  Guida Utente — JDER (Java Diagrammi E/R)
%  Compilare con: pdflatex GuidaJDER.tex  (due passaggi)
%  Dipendenze: geometry, hyperref, graphicx, booktabs,
%              tcolorbox, fontenc, inputenc, babel,
%              enumitem, xcolor, titlesec, parskip,
%              microtype, fancyhdr, lastpage
% ============================================================
\documentclass[a4paper,11pt]{article}

% --- Encoding & lingua ----------------------------------------
\usepackage[T1]{fontenc}
\usepackage[utf8]{inputenc}
\usepackage[italian]{babel}

% --- Layout ---------------------------------------------------
\usepackage[
  top=2.5cm, bottom=2.8cm,
  left=2.8cm, right=2.8cm,
  headheight=14pt
]{geometry}
\usepackage{parskip}         % spazio tra paragrafi al posto del rientro
\usepackage{microtype}       % micro-tipografia
\usepackage{emptypage}       % pagine vuote realmente vuote

% --- Colori ---------------------------------------------------
\usepackage{xcolor}
% Palette brand JDER (dal sito jder.it)
\definecolor{jderorange}{RGB}{240,120,24}   % primario sito
\definecolor{jderred}{RGB}{233,75,60}        % accento sito
\definecolor{jderdarkblue}{RGB}{30,60,114}   % heading h2/h3 sito
\definecolor{jderblue}{RGB}{42,82,152}       % link/badge sito
\definecolor{jderdark}{RGB}{44,62,80}        % testo body sito
\definecolor{jdergray}{RGB}{80,90,100}       % testo secondario
\definecolor{coverbgA}{RGB}{15,32,39}        % inizio gradiente sfondo sito
\definecolor{coverbgB}{RGB}{32,58,67}        % metà gradiente sito
\definecolor{notebg}{RGB}{255,248,235}       % caldo-arancio chiaro
\definecolor{warnbg}{RGB}{255,235,230}       % rosso chiaro
\definecolor{tipbg}{RGB}{235,245,255}        % blu chiaro

% --- Tipografia -----------------------------------------------
\usepackage{lmodern}
\usepackage{textcomp}

% --- Grafica --------------------------------------------------
\usepackage{graphicx}
\usepackage{float}

% --- Tabelle --------------------------------------------------
\usepackage{booktabs}
\usepackage{array}
\usepackage{longtable}

% --- Box colorati ---------------------------------------------
\usepackage[most]{tcolorbox}

\tcbset{
  notebox/.style={
    colback=notebg, colframe=jderorange,
    coltitle=jderdarkblue,
    fonttitle=\bfseries, left=6pt, right=6pt,
    top=4pt, bottom=4pt, arc=4pt
  },
  warnbox/.style={
    colback=warnbg, colframe=jderred
    
    ,
    coltitle=jderred,
    fonttitle=\bfseries, left=6pt, right=6pt,
    top=4pt, bottom=4pt, arc=4pt
  },
  tipbox/.style={
    colback=tipbg, colframe=jderblue,
    coltitle=jderdarkblue,
    fonttitle=\bfseries, left=6pt, right=6pt,
    top=4pt, bottom=4pt, arc=4pt
  },
  shortcutbox/.style={
    colback=gray!8, colframe=jderorange!40,
    left=6pt, right=6pt, top=4pt, bottom=4pt, arc=4pt
  }
}

% --- Elenchi puntati ------------------------------------------
\usepackage{enumitem}
\setlist[itemize]{leftmargin=1.6em, itemsep=2pt, topsep=4pt}
\setlist[enumerate]{leftmargin=1.6em, itemsep=2pt, topsep=4pt}

% --- Titoli di sezione ----------------------------------------
\usepackage{titlesec}
\titleformat{\section}
  {\Large\bfseries\color{jderorange}}
  {\thesection}{1em}{}[\vspace{-4pt}\textcolor{jderorange}{\rule{\linewidth}{1.5pt}}]
\titleformat{\subsection}
  {\large\bfseries\color{jderdarkblue}}
  {\thesubsection}{0.8em}{}
\titleformat{\subsubsection}
  {\normalsize\bfseries\color{jderblue}}
  {\thesubsubsection}{0.7em}{}

% --- Intestazioni e piè di pagina -----------------------------
\usepackage{fancyhdr}
\usepackage{lastpage}
\pagestyle{fancy}
\fancyhf{}
\fancyhead[L]{\small\color{jderdarkblue}\textbf{JDER} --- Guida Utente}
\fancyhead[R]{\small\color{jdergray}\leftmark}
\fancyfoot[C]{\small\color{jdergray}Pagina \thepage\ di \pageref{LastPage}}
\renewcommand{\headrulewidth}{0.4pt}
\renewcommand{\footrulewidth}{0.3pt}

% --- Hyperlink ------------------------------------------------
\usepackage[
  colorlinks=true,
  linkcolor=jderdarkblue,
  urlcolor=jderorange,
  pdfauthor={JDER Project},
  pdftitle={Guida Utente JDER},
  pdfsubject={Editor di Diagrammi UML e E/R}
]{hyperref}

% --- Comandi personalizzati -----------------------------------
\newcommand{\key}[1]{\texttt{\textbf{#1}}}
\newcommand{\menu}[1]{\textit{#1}}
\newcommand{\uielement}[1]{\textbf{#1}}
\newcommand{\jder}{\textbf{JDER}}
\newcommand{\theauthor}{LoryPelli}

% =============================================================
\begin{document}


    
    % ===================== COPERTINA =============================
\begin{titlepage}
  \pagecolor{coverbgA}
  \color{white}
  \centering
  \vspace*{2.5cm}

  % Decorazione superiore
  {\color{jderorange}\rule{4cm}{2pt}}\quad
  {\color{white!50!coverbgA}\rule{2cm}{0.5pt}}\quad
  {\color{jderred}\rule{4cm}{2pt}}

  \vspace{1.2cm}

  % Titolo principale in arancione brand
  {\fontsize{60}{68}\selectfont\bfseries\color{jderorange} JDER}

  \vspace{0.3cm}
  {\large\color{white!80!coverbgA} Java Diagrammi E/R}

  \vspace{0.8cm}
  {\color{jderorange}\rule{9cm}{1pt}}

  \vspace{0.8cm}
  {\LARGE\bfseries Guida Utente Completa}

  \vspace{0.5cm}
  {\normalsize\color{white!70!coverbgA} a cura di \textbf{\color{jderorange}\theauthor}}

  \vspace{2cm}
  {\normalsize\color{white!85!coverbgA}
    Editor moderno e multipiattaforma per\\
    \textbf{Diagrammi Entità-Relazione},
    \textbf{Diagrammi delle Classi}\\
    e \textbf{Diagrammi dei Casi d'Uso}
  }

  \vspace{1.5cm}

  % Badge tecnologie
  \tcbox[on line, colback=coverbgB, colframe=jderorange,
    arc=4pt, boxrule=0.8pt, left=4pt, right=4pt, top=3pt, bottom=3pt,
    fontupper=\bfseries\color{white}]{Kotlin}\quad
  \tcbox[on line, colback=coverbgB, colframe=jderorange,
    arc=4pt, boxrule=0.8pt, left=4pt, right=4pt, top=3pt, bottom=3pt,
    fontupper=\bfseries\color{white}]{Compose Multiplatform}\quad
  \tcbox[on line, colback=coverbgB, colframe=jderorange,
    arc=4pt, boxrule=0.8pt, left=4pt, right=4pt, top=3pt, bottom=3pt,
    fontupper=\bfseries\color{white}]{Material 3}

  \vfill
  {\color{jderorange}\rule{9cm}{0.8pt}}\\[0.4cm]
  {\small\color{white!70!coverbgA} Basato su Kotlin \& Jetpack Compose Multiplatform \quad\textbullet\quad Licenza MIT}\\[0.3cm]
  {\small\color{jderorange} \url{https://jder.it}}
\end{titlepage}
\nopagecolor

% ===================== INDICE ================================
\tableofcontents
\newpage

% =============================================================
\section{Introduzione}
\label{sec:intro}

\jder{} (\textit{Java Diagrammi E/R}) è un editor grafico open-source per la creazione di diagrammi
di modellazione software e dati. L'applicazione è una riscrittura completa in \textbf{Kotlin}
con \textbf{Jetpack Compose Multiplatform}, che sostituisce e migliora le versioni precedenti
sviluppate da Alessandro Ballini e successivamente mantenute da Gianvito Pio.

\subsection{Cosa si può fare con JDER}

\jder{} supporta tre tipologie di diagramma distinte, ognuna accessibile tramite una scheda
dedicata nell'interfaccia principale:

\begin{itemize}
  \item \textbf{Diagramma Entità-Relazione (E/R)} --- per modellare la struttura logica di una
        base di dati secondo la notazione classica con entità, relazioni e attributi.
  \item \textbf{Diagramma delle Classi UML} --- per descrivere la struttura statica di un sistema
        orientato agli oggetti con classi, interfacce, enumerazioni e le loro relazioni.
  \item \textbf{Diagramma dei Casi d'Uso UML} --- per rappresentare le interazioni tra attori
        esterni e i casi d'uso di un sistema software.
\end{itemize}

\subsection{Requisiti di sistema}

\begin{itemize}
  \item Sistema operativo: Windows, macOS o Linux con Java 17 o superiore installato.
  \item RAM consigliata: almeno 512\,MB liberi.
  \item Risoluzione schermo: minima 800$\times$600 pixel (si consiglia Full HD o superiore).
\end{itemize}

\subsection{Avvio dell'applicazione}

Scaricare il file \texttt{JDER.jar} dalla pagina delle release su GitHub
(\url{https://github.com/LoryPelli/JDER/releases/latest}) oppure dal sito ufficiale
(\url{https://jder.it}).

Per avviare l'applicazione:
\begin{enumerate}
  \item Doppio clic sul file \texttt{JDER.jar} (se Java è associato all'estensione \texttt{.jar}).
  \item Oppure, da terminale: \texttt{java -jar JDER.jar}
\end{enumerate}

\begin{tcolorbox}[warnbox, title={Attenzione}]
  Se il doppio clic non avvia l'applicazione, verificare che Java 17 o superiore sia installato
  e che la variabile di ambiente \texttt{JAVA\_HOME} sia configurata correttamente.
\end{tcolorbox}

% =============================================================
\section{Interfaccia Principale}
\label{sec:interfaccia}

All'avvio, \jder{} mostra una finestra massimizzata con i seguenti componenti principali:

\begin{enumerate}
  \item \textbf{Barra delle schede} (in alto) --- permette di passare tra i tre tipi di diagramma.
  \item \textbf{Barra degli strumenti} --- contiene i menu e i pulsanti di azione.
  \item \textbf{Barra degli strumenti di disegno} --- sotto la barra principale, contiene gli
        strumenti di creazione degli elementi.
  \item \textbf{Lavagna (Canvas)} --- area centrale dove vengono disegnati i diagrammi.
  \item \textbf{Pannello Proprietà} --- barra laterale destra che appare quando si seleziona un
        elemento, mostrando e permettendo di modificare le sue proprietà.
\end{enumerate}

\subsection{Barra delle schede}

La barra delle schede in cima all'interfaccia presenta tre schede:

\begin{itemize}
  \item \uielement{Diagramma E/R} --- Diagramma Entità-Relazione
  \item \uielement{Diagramma Classi} --- Diagramma delle Classi
  \item \uielement{Casi d'Uso} --- Diagramma dei Casi d'Uso
\end{itemize}

Ogni scheda mantiene il proprio stato indipendente: è possibile lavorare su un diagramma E/R e
contemporaneamente tenere aperto un diagramma delle classi in un'altra scheda, passando da una
all'altra senza perdere le modifiche non salvate.

\begin{tcolorbox}[notebox, title={Nota}]
  Alla chiusura dell'applicazione, se ci sono modifiche non salvate in \emph{qualsiasi} scheda,
  \jder{} mostrerà una finestra di conferma prima di uscire.
\end{tcolorbox}

\subsection{Barra degli strumenti principale}

La barra degli strumenti principale (TopAppBar) contiene, da sinistra a destra:

\begin{itemize}
  \item \textbf{Titolo} --- mostra il nome del file aperto e un asterisco (\texttt{*}) se ci sono
        modifiche non salvate.
  \item Menu \uielement{File} --- gestione di file (nuovo, apri, salva, salva con nome, elimina).
  \item Menu \uielement{Vista} --- controlli di zoom (Zoom In, Zoom Out, Reimposta Zoom).
  \item Menu \uielement{Esporta} --- esportazione del diagramma come immagine PNG.
  \item \textbf{Selettore palette} --- un cerchio colorato che apre un selettore per scegliere
        tra le 20 palette di colori Material disponibili.
  \item \textbf{Interruttore tema} --- alterna tra modalità chiara e modalità scura.
\end{itemize}

\subsection{Barra degli strumenti di disegno}

Sotto la barra principale si trova la barra degli strumenti di disegno, che varia in base alla
scheda attiva (vedi sezioni successive). Contiene sempre:

\begin{itemize}
  \item Pulsante \uielement{Seleziona e Sposta} (freccia): modalità di selezione e spostamento.
  \item Pulsanti \uielement{Annulla} e \uielement{Ripristina}.
  \item Strumenti di creazione specifici per il tipo di diagramma.
  \item Pulsante \uielement{Salva} rapido.
  \item Pulsante \uielement{Elimina} (attivo solo con un elemento selezionato).
  \item Indicatore e controlli di zoom (a destra).
\end{itemize}

\subsection{La lavagna}

La lavagna è l'area principale di disegno. Supporta le seguenti interazioni:

\begin{center}
\begin{tabular}{@{}ll@{}}
  \toprule
  \textbf{Azione} & \textbf{Controllo} \\
  \midrule
  Zoom avanti & \key{Ctrl} + \key{+} oppure scroll del mouse verso l'alto \\
  Zoom indietro & \key{Ctrl} + \key{-} oppure scroll del mouse verso il basso \\
  Scorrimento (pan) & Trascinamento con il tasto centrale del mouse \\
  Selezionare un elemento & Clic sinistro sull'elemento \\
  Spostare un elemento & Clic sinistro + trascinamento (in modalità Seleziona) \\
  Menu contestuale & Clic destro sull'elemento \\
  Creare un elemento & Clic sinistro sull'area vuota (con strumento attivo) \\
  Deselezionare / Annullare & \key{Esc} \\
  \bottomrule
\end{tabular}
\end{center}

\begin{tcolorbox}[tipbox, title={Suggerimento}]
  Lo zoom massimo è al \textbf{300\%} e il minimo al \textbf{25\%}. Per tornare allo zoom
  predefinito (100\%) usare \menu{Vista} $\to$ \menu{Reimposta Zoom} oppure \key{Ctrl} + \key{0}.
\end{tcolorbox}

\subsection{Pannello Proprietà}

Il Pannello Proprietà si apre automaticamente sul lato destro della lavagna quando si seleziona
un elemento. Mostra le informazioni dell'elemento selezionato e permette di:

\begin{itemize}
  \item Modificare le proprietà principali (nome, tipo, documentazione, ecc.).
  \item Aggiungere, modificare o eliminare attributi/membri dell'elemento.
  \item Aggiungere o modificare connessioni (nel diagramma E/R).
  \item Eseguire operazioni speciali (es.\@ conversione in entità associativa).
\end{itemize}

Il pannello si chiude tramite il pulsante \textbf{X} in alto a destra o deselezionando
l'elemento.

% =============================================================
\section{Diagramma Entità-Relazione (E/R)}
\label{sec:er}

Il modulo E/R implementa la notazione classica dei diagrammi Entità-Relazione, ampiamente
utilizzata nella progettazione concettuale di basi di dati.

\subsection{Strumenti disponibili}

La barra degli strumenti del diagramma E/R offre i seguenti strumenti di creazione:

\begin{itemize}
  \item \uielement{Seleziona e Sposta} (icona freccia) --- seleziona e sposta elementi.
  \item \uielement{Crea Entità} (icona rettangolo) --- crea una nuova entità sulla lavagna.
  \item \uielement{Crea Relazione} (icona rombo) --- crea una nuova relazione sulla lavagna.
  \item \uielement{Crea Nota} (icona foglietto) --- crea una nota testuale.
\end{itemize}

\subsection{Entità}

Un'entità rappresenta una categoria di oggetti del dominio applicativo. In \jder{} si distinguono
due tipi di entità:

\subsubsection{Entità forte}

È il tipo predefinito. Viene rappresentata graficamente con un rettangolo con bordi arrotondati
e il nome dell'entità al centro.

\textbf{Creazione:}
\begin{enumerate}
  \item Selezionare lo strumento \uielement{Crea Entità}.
  \item Fare clic nell'area vuota della lavagna nel punto desiderato.
\end{enumerate}

\textbf{Proprietà modificabili} (tramite menu contestuale o pannello proprietà):
\begin{itemize}
  \item \textbf{Nome} --- il nome dell'entità.
  \item \textbf{Tipo} --- Forte (predefinito) o Debole.
  \item \textbf{Documentazione} --- testo libero di descrizione.
\end{itemize}

\subsubsection{Entità debole}

Un'entità debole è un'entità la cui identificazione dipende da un'altra entità. Viene
rappresentata con un doppio rettangolo (bordo interno ed esterno).

Per creare un'entità debole, editare le proprietà di un'entità esistente e attivare l'opzione
\uielement{Tipo: Debole} nella finestra di dialogo.

\subsection{Relazioni}

Una relazione rappresenta un legame semantico tra due o più entità. Viene rappresentata
graficamente con un rombo con il nome al centro.

\textbf{Creazione:}
\begin{enumerate}
  \item Selezionare lo strumento \uielement{Crea Relazione}.
  \item Fare clic nell'area vuota della lavagna.
\end{enumerate}

\textbf{Proprietà modificabili:}
\begin{itemize}
  \item \textbf{Nome} --- il nome della relazione.
  \item \textbf{Documentazione} --- descrizione testuale.
\end{itemize}

\subsubsection{Aggiungere connessioni a una relazione}

Dopo aver creato una relazione, occorre collegarla alle entità che partecipano. Per aggiungere
una connessione:

\begin{enumerate}
  \item Selezionare la relazione sulla lavagna.
  \item Nel pannello Proprietà, fare clic su \uielement{Aggiungi Connessione}, oppure usare il
        menu contestuale (clic destro) $\to$ \uielement{Aggiungi Connessione}.
  \item Nella finestra di dialogo, selezionare l'entità da collegare e la cardinalità.
\end{enumerate}

\subsubsection{Cardinalità supportate}

\jder{} supporta le seguenti notazioni di cardinalità:

\begin{center}
\begin{tabular}{@{}lll@{}}
  \toprule
  \textbf{Valore} & \textbf{Simbolo} & \textbf{Significato} \\
  \midrule
  \texttt{1}      & \texttt{1}       & Esattamente uno \\
  \texttt{(0,1)}  & \texttt{(0,1)}   & Zero o uno \\
  \texttt{(1,1)}  & \texttt{(1,1)}   & Esattamente uno (con vincolo) \\
  \texttt{N}      & \texttt{N}       & Molti (senza limite) \\
  \texttt{(0,N)}  & \texttt{(0,N)}   & Zero o molti \\
  \texttt{(1,N)}  & \texttt{(1,N)}   & Uno o molti \\
  \bottomrule
\end{tabular}
\end{center}

\subsubsection{Conversione in entità associativa}

Quando una relazione coinvolge due entità con cardinalità N-a-N (entrambe N, (0,N) o (1,N)),
è possibile convertirla automaticamente in un'\textbf{entità associativa}:

\begin{enumerate}
  \item Selezionare la relazione N-a-N.
  \item Fare clic destro $\to$ \uielement{Converti in Entità Associativa}, oppure usare il
        pannello proprietà.
\end{enumerate}

Il risultato è un'entità con lo stesso nome della relazione, collegata alle due entità originali
con le cardinalità appropriate. Questa operazione è disponibile solo per relazioni con esattamente
due connessioni entrambe di tipo ``molti''.

\subsection{Attributi}

Gli attributi descrivono le proprietà di un'entità o di una relazione. \jder{} supporta cinque
tipi di attributo:

\begin{center}
\begin{tabular}{@{}lll@{}}
  \toprule
  \textbf{Tipo} & \textbf{Rappresentazione grafica} & \textbf{Descrizione} \\
  \midrule
  Normale    & Ellisse con bordo semplice & Attributo semplice \\
  Chiave     & Ellisse con testo sottolineato & Identificatore primario \\
  Multivalore & Ellisse con doppio bordo & Può avere più valori \\
  Derivato   & Ellisse con bordo tratteggiato & Calcolato da altri attributi \\
  Composto   & Ellisse + sotto-attributi collegati & Composto da più parti \\
  \bottomrule
\end{tabular}
\end{center}

\textbf{Aggiungere un attributo:}
\begin{enumerate}
  \item Selezionare un'entità o una relazione.
  \item Nel pannello Proprietà, fare clic su \uielement{Aggiungi Attributo}, oppure usare il
        menu contestuale $\to$ \uielement{Aggiungi Attributo}.
  \item Compilare il dialogo:
        \begin{itemize}
          \item \textbf{Nome} --- obbligatorio.
          \item \textbf{Chiave Primaria} --- attivare se l'attributo identifica l'entità.
          \item \textbf{Tipo} --- scegliere tra Normale, Chiave, Multivalore, Derivato, Composto.
          \item \textbf{Molteplicità} --- visibile solo per gli attributi Multivalore
                (es.\@ \texttt{1..N}, \texttt{0..N}).
          \item \textbf{Componenti} --- visibile solo per gli attributi Composti; è possibile
                aggiungere sotto-attributi con il pulsante \uielement{+ Aggiungi Componente}.
        \end{itemize}
\end{enumerate}

\begin{tcolorbox}[notebox, title={Nota}]
  Selezionando il tipo \textbf{Chiave}, l'opzione ``Chiave Primaria'' viene attivata
  automaticamente e non è modificabile manualmente.
\end{tcolorbox}

\subsection{Note testuali}

Le note consentono di aggiungere annotazioni libere alla lavagna.

\textbf{Creazione:}
\begin{enumerate}
  \item Selezionare lo strumento \uielement{Crea Nota}.
  \item Fare clic nella posizione desiderata.
\end{enumerate}

Il testo della nota si modifica tramite il menu contestuale (clic destro) oppure
selezionando la nota e usando il pannello Proprietà.

\subsection{Menu contestuale nel diagramma E/R}

Il clic destro su un elemento della lavagna apre un menu contestuale con azioni rapide:

\begin{itemize}
  \item \textbf{Su un'entità:} Modifica Proprietà, Aggiungi Attributo, Elimina.
  \item \textbf{Su una relazione:} Modifica Proprietà, Aggiungi Attributo, Aggiungi Connessione,
        Converti in Entità Associativa (abilitato solo per relazioni N-a-N), Elimina.
  \item \textbf{Su una nota:} Modifica Testo, Elimina.
\end{itemize}

% =============================================================
\section{Diagramma delle Classi}
\label{sec:class}

Il modulo Diagramma delle Classi implementa la notazione UML per rappresentare la struttura
statica di un sistema a oggetti.

\subsection{Strumenti disponibili}

\begin{itemize}
  \item \uielement{Seleziona e Sposta} --- seleziona e sposta elementi.
  \item \uielement{Crea Classe} (icona rettangolo) --- crea una nuova classe.
  \item \uielement{Crea Interfaccia} (icona ellisse) --- crea una nuova interfaccia.
  \item \uielement{Crea Enum} (icona rombo) --- crea una nuova enumerazione.
  \item \uielement{Crea Relazione} (icona freccia) --- disegna una relazione tra due elementi.
  \item \uielement{Crea Nota} (icona foglietto) --- crea una nota testuale.
\end{itemize}

\subsection{Elementi del diagramma delle classi}

\subsubsection{Classe}

La classe è il blocco fondamentale. Viene rappresentata con un rettangolo diviso in tre sezioni:
nome (in alto), attributi (al centro) e metodi (in basso).

\textbf{Tipi di classe supportati:}
\begin{center}
\begin{tabular}{@{}lll@{}}
  \toprule
  \textbf{Tipo} & \textbf{Stereotipo UML} & \textbf{Note} \\
  \midrule
  Classe          & (nessuno)         & Tipo predefinito \\
  Interfaccia     & \texttt{<<interface>>}   & \\
  Enumerazione    & \texttt{<<enum>>}        & \\
  Classe Astratta & \texttt{<<abstract>>}    & Il nome è in corsivo \\
  \bottomrule
\end{tabular}
\end{center}

\textbf{Creazione:}
\begin{enumerate}
  \item Selezionare lo strumento desiderato (\uielement{Crea Classe}, \uielement{Crea Interfaccia}
        o \uielement{Crea Enum}).
  \item Fare clic sull'area vuota della lavagna.
\end{enumerate}

\textbf{Proprietà della classe} (dialogo di modifica):
\begin{itemize}
  \item \textbf{Nome} --- obbligatorio.
  \item \textbf{Tipo} --- Classe, Interfaccia, Enum, Classe Astratta (con selezione a radio).
  \item \textbf{Astratta} --- disponibile solo per le classi; rende il nome in corsivo.
  \item \textbf{Documentazione} --- campo di testo libero.
\end{itemize}

\subsubsection{Attributi e metodi di una classe}

Ogni classe può avere un numero arbitrario di attributi e di metodi.

\textbf{Aggiungere un attributo:}
\begin{enumerate}
  \item Selezionare la classe.
  \item Nel pannello Proprietà, fare clic su \uielement{Aggiungi Attributo}.
\end{enumerate}

\textbf{Aggiungere un metodo:}
\begin{enumerate}
  \item Selezionare la classe.
  \item Nel pannello Proprietà, fare clic su \uielement{Aggiungi Metodo}.
\end{enumerate}

\textbf{Campi del dialogo per attributi e metodi:}
\begin{itemize}
  \item \textbf{Nome} --- obbligatorio.
  \item \textbf{Tipo} --- tipo dell'attributo o tipo di ritorno del metodo.
  \item \textbf{Visibilità} --- Pubblico (\texttt{+}), Privato (\texttt{-}),
        Protetto (\texttt{\#}), Package (\texttt{\textasciitilde}).
  \item \textbf{Statico} --- se attivato, il nome viene sottolineato nella rappresentazione.
  \item \textbf{Astratto} --- disponibile solo per i metodi; il nome viene reso in corsivo.
\end{itemize}

\subsection{Relazioni tra classi}

\jder{} supporta i sei tipi di relazione UML standard:

\begin{center}
\begin{tabular}{@{}lll@{}}
  \toprule
  \textbf{Tipo} & \textbf{Notazione} & \textbf{Descrizione} \\
  \midrule
  Associazione  & Linea con freccia aperta & Dipendenza semantica generica \\
  Aggregazione  & Linea con rombo vuoto    & Parte di (``ha un'') \\
  Composizione  & Linea con rombo pieno    & Parte di (con dipendenza di vita) \\
  Ereditarietà  & Linea con triangolo vuoto & Generalizzazione (``è un'') \\
  Realizzazione & Linea tratteggiata con triangolo & Implementazione interfaccia \\
  Dipendenza    & Linea tratteggiata con freccia & Dipendenza d'uso \\
  \bottomrule
\end{tabular}
\end{center}

\textbf{Creazione di una relazione:}
\begin{enumerate}
  \item Selezionare lo strumento \uielement{Crea Relazione}.
  \item Fare clic sulla classe \textbf{sorgente}.
  \item Fare clic sulla classe \textbf{destinazione}.
  \item Nella finestra di dialogo, selezionare il tipo di relazione e compilare
        i campi opzionali: etichetta, molteplicità sorgente e molteplicità destinazione.
\end{enumerate}

\textbf{Campi della relazione:}
\begin{itemize}
  \item \textbf{Tipo} --- uno dei sei tipi elencati.
  \item \textbf{Etichetta} (opzionale) --- testo visibile sulla linea della relazione.
  \item \textbf{Molteplicità sorgente} (opzionale) --- es.\@ \texttt{1}, \texttt{0..*}.
  \item \textbf{Molteplicità destinazione} (opzionale) --- es.\@ \texttt{1..*}.
\end{itemize}

% =============================================================
\section{Diagramma dei Casi d'Uso}
\label{sec:usecase}

Il modulo Casi d'Uso implementa la notazione UML per i diagrammi dei casi d'uso, utili per
descrivere i requisiti funzionali di un sistema dal punto di vista degli attori esterni.

\subsection{Strumenti disponibili}

\begin{itemize}
  \item \uielement{Seleziona e Sposta} --- seleziona e sposta elementi.
  \item \uielement{Crea Attore} (icona omino) --- crea un attore.
  \item \uielement{Crea Caso d'Uso} (icona ellisse) --- crea un caso d'uso.
  \item \uielement{Crea Relazione} (icona freccia) --- disegna una relazione.
  \item \uielement{Crea Nota} (icona foglietto) --- crea una nota testuale.
  \item \uielement{Crea Sistema} (icona rettangolo vuoto) --- crea un confine di sistema.
\end{itemize}

\subsection{Attori}

Un attore rappresenta un'entità esterna al sistema (utente, sistema esterno, timer, ecc.)
che interagisce con esso. Viene rappresentato come un omino stilizzato (``stick figure'') con il
nome sotto.

\textbf{Creazione:}
\begin{enumerate}
  \item Selezionare lo strumento \uielement{Crea Attore}.
  \item Fare clic sull'area vuota.
\end{enumerate}

\textbf{Modifica:} clic sull'attore per selezionarlo, poi usare il menu contestuale
o il pannello proprietà per cambiare il nome.

\subsection{Casi d'Uso}

Un caso d'uso rappresenta una funzionalità o un servizio offerto dal sistema. Viene rappresentato
con un'ellisse contenente il nome.

\textbf{Creazione:}
\begin{enumerate}
  \item Selezionare lo strumento \uielement{Crea Caso d'Uso}.
  \item Fare clic sull'area vuota.
\end{enumerate}

\textbf{Proprietà:}
\begin{itemize}
  \item \textbf{Nome} --- obbligatorio.
  \item \textbf{Documentazione} --- descrizione testuale del caso d'uso.
\end{itemize}

\subsection{Relazioni tra elementi}

\jder{} supporta quattro tipi di relazione per i diagrammi dei casi d'uso:

\begin{center}
\begin{tabular}{@{}lll@{}}
  \toprule
  \textbf{Tipo} & \textbf{Notazione} & \textbf{Utilizzo tipico} \\
  \midrule
  Associazione     & Linea semplice                  & Attore $\leftrightarrow$ Caso d'uso \\
  Include          & Freccia tratteggiata \texttt{<<include>>} & Inclusione obbligatoria \\
  Extend           & Freccia tratteggiata \texttt{<<extend>>}  & Estensione condizionale \\
  Generalizzazione & Linea con triangolo vuoto        & Specializzazione \\
  \bottomrule
\end{tabular}
\end{center}

\textbf{Creazione di una relazione:}
\begin{enumerate}
  \item Selezionare lo strumento \uielement{Crea Relazione}.
  \item Fare clic sull'elemento \textbf{sorgente}.
  \item Fare clic sull'elemento \textbf{destinazione}.
  \item Selezionare il tipo di relazione nella finestra di dialogo.
\end{enumerate}

\subsection{Confini di sistema}

Il confine di sistema (\textit{System Boundary}) è un rettangolo che raggruppa visivamente i
casi d'uso appartenenti allo stesso sistema, separandoli dagli attori esterni.

\textbf{Creazione:}
\begin{enumerate}
  \item Selezionare lo strumento \uielement{Crea Sistema}.
  \item Fare clic sull'area vuota.
  \item Il rettangolo creato può essere ridimensionato trascinandone i bordi.
\end{enumerate}

\textbf{Proprietà:} nome del sistema (visualizzato in cima al rettangolo).

% =============================================================
\section{Operazioni Comuni}
\label{sec:comuni}

Questa sezione descrive le funzionalità condivise tra tutti e tre i tipi di diagramma.

\subsection{Gestione dei file}

\subsubsection{Nuovo diagramma}

Per creare un nuovo diagramma vuoto:
\begin{itemize}
  \item Menu \menu{File} $\to$ \menu{Nuovo Diagramma}
  \item Scorciatoia: \key{Ctrl} + \key{N}
\end{itemize}

Se esistono modifiche non salvate, \jder{} chiede conferma prima di procedere.

\subsubsection{Aprire un diagramma}

Per aprire un file di diagramma esistente:
\begin{itemize}
  \item Menu \menu{File} $\to$ \menu{Apri Diagramma\ldots}
  \item Scorciatoia: \key{Ctrl} + \key{O}
\end{itemize}

Si aprirà il \textbf{Gestore File} integrato (vedi sezione~\ref{sec:filemanager}).

\subsubsection{Salvare un diagramma}

\textbf{Salva} (sovrascrive il file corrente):
\begin{itemize}
  \item Menu \menu{File} $\to$ \menu{Salva}
  \item Pulsante \uielement{Salva} nella barra di disegno
  \item Scorciatoia: \key{Ctrl} + \key{S}
\end{itemize}

\textbf{Salva con nome} (sceglie una nuova posizione o nome):
\begin{itemize}
  \item Menu \menu{File} $\to$ \menu{Salva con nome\ldots}
\end{itemize}

I diagrammi vengono salvati in formato \textbf{JSON}. Tutti i file utilizzano l'estensione \texttt{.json},
indipendentemente dal tipo di diagramma.

\begin{tcolorbox}[notebox, title={Nota}]
  Il titolo nella barra degli strumenti mostra un asterisco (\texttt{*}) dopo il nome del file
  per indicare la presenza di modifiche non salvate.
\end{tcolorbox}

\subsection{Il Gestore File}
\label{sec:filemanager}

\jder{} include un gestore file integrato che funziona sia in modalità \textbf{Apri} che in
modalità \textbf{Salva con nome}. Le sue funzionalità sono:

\begin{itemize}
  \item \textbf{Navigazione delle directory:} fare clic sulle cartelle per entrarvi; usare il
        pulsante Indietro per tornare alla directory superiore.
  \item \textbf{Barra del percorso:} mostra e permette di modificare il percorso corrente.
        Premere \key{Invio} per navigare a un percorso digitato manualmente.
  \item \textbf{Filtro file:} attivare lo switch \uielement{Mostra solo file supportati} per
        visualizzare solo i file con l'estensione corretta.
  \item \textbf{Aggiornamento:} pulsante di aggiornamento per ricaricare il contenuto della
        directory.
  \item \textbf{Conferma sovrascrittura:} se si salva con il nome di un file esistente,
        \jder{} chiede conferma prima di sovrascriverlo.
\end{itemize}

\subsection{Esportazione come PNG}

È possibile esportare qualsiasi diagramma come immagine PNG ad alta qualità:

\begin{enumerate}
  \item Menu \menu{Esporta} $\to$ \menu{Esporta come PNG\ldots}
  \item Scegliere la cartella di destinazione e il nome del file nel Gestore File.
\end{enumerate}

L'immagine esportata cattura l'intero contenuto della lavagna, inclusi tutti gli elementi
(anche quelli non visibili nell'area di viewport corrente).

\begin{tcolorbox}[tipbox, title={Suggerimento}]
  Per ottenere un'immagine di dimensioni adeguate, assicurarsi che tutti gli elementi del
  diagramma siano posizionati in modo ordinato prima di esportare.
\end{tcolorbox}

\subsection{Annulla e Ripristina}

\jder{} mantiene una cronologia completa delle azioni per ogni diagramma, permettendo di
annullare e ripristinare qualsiasi operazione:

\begin{itemize}
  \item \textbf{Annulla:} \key{Ctrl} + \key{Z} oppure pulsante \uielement{Annulla} (icona freccia sinistra) nella barra.
  \item \textbf{Ripristina:} \key{Ctrl} + \key{Y} oppure pulsante \uielement{Ripristina} (icona freccia destra) nella barra.
\end{itemize}

I pulsanti sono disabilitati quando non ci sono azioni da annullare o ripristinare.

\subsection{Eliminazione di elementi}

Per eliminare un elemento selezionato:
\begin{itemize}
  \item Tasto \key{Delete} o \key{Backspace}.
  \item Pulsante \uielement{Elimina} (icona cestino) nella barra degli strumenti.
  \item Menu \menu{File} $\to$ \menu{Elimina Elemento Selezionato}.
  \item Menu contestuale $\to$ \uielement{Elimina}.
\end{itemize}

\begin{tcolorbox}[warnbox, title={Attenzione}]
  L'eliminazione è un'operazione reversibile tramite \textbf{Annulla} (\key{Ctrl} + \key{Z}).
  Tuttavia, l'eliminazione di un'entità nel diagramma E/R elimina anche tutte le connessioni
  che la coinvolgono nelle relazioni associate.
\end{tcolorbox}

\subsection{Zoom e navigazione}

I controlli di zoom sono disponibili tramite diversi metodi:

\begin{center}
\begin{tabular}{@{}lll@{}}
  \toprule
  \textbf{Azione} & \textbf{Scorciatoia} & \textbf{Menu} \\
  \midrule
  Zoom In         & \key{Ctrl} + \key{+} & Vista $\to$ Zoom In \\
  Zoom Out        & \key{Ctrl} + \key{-} & Vista $\to$ Zoom Out \\
  Reimposta Zoom  & \key{Ctrl} + \key{0} & Vista $\to$ Reimposta Zoom \\
  Scroll su/giù  & Rotella del mouse    & --- \\
  \bottomrule
\end{tabular}
\end{center}

Lo zoom è visualizzato come percentuale nella barra degli strumenti (es.\@ \texttt{100\%}).
I limiti sono: minimo \textbf{25\%}, massimo \textbf{300\%}.

\subsection{Personalizzazione visiva}

\subsubsection{Modalità chiara e scura}

\jder{} supporta due modalità di visualizzazione:
\begin{itemize}
  \item \textbf{Modalità chiara} (Light Mode) --- sfondo bianco, testo scuro.
  \item \textbf{Modalità scura} (Dark Mode) --- sfondo scuro, testo chiaro.
\end{itemize}

Per cambiare modalità, fare clic sul pulsante con l'icona sole/luna nella barra degli strumenti.
La transizione è animata in modo fluido.

\subsubsection{Palette di colori}

\jder{} offre 20 palette di colori Material Design ufficiali selezionabili a runtime.
Per cambiarla, fare clic sul \textbf{cerchio colorato} nella barra degli strumenti e selezionare
la palette desiderata dal selettore a griglia.

Le palette disponibili sono:
\begin{itemize}
  \item Rosso, Rosa, Viola, Viola Scuro, Indaco, Blu, Blu Chiaro, Ciano,
        Verde Acqua, Verde, Verde Chiaro, Lime, Giallo, Ambra, Arancione,
        Arancione Scuro, Marrone, Grigio, Grigio Blu, Bianco \& Nero.
\end{itemize}

La palette selezionata determina il colore principale (primary) dell'intera interfaccia,
inclusi pulsanti, indicatori e accenti.

% =============================================================
\section{Scorciatoie da Tastiera}
\label{sec:shortcuts}

La tabella seguente riassume tutte le scorciatoie da tastiera disponibili in \jder{}:

\begin{center}
\begin{tabular}{@{}ll@{}}
  \toprule
  \textbf{Scorciatoia} & \textbf{Azione} \\
  \midrule
  \key{Ctrl} + \key{N}           & Nuovo Diagramma \\
  \key{Ctrl} + \key{O}           & Apri Diagramma \\
  \key{Ctrl} + \key{S}           & Salva (sovrascrive il file corrente) \\
  \key{Ctrl} + \key{Z}           & Annulla \\
  \key{Ctrl} + \key{Y}           & Ripristina \\
  \key{Ctrl} + \key{+}           & Zoom In \\
  \key{Ctrl} + \key{-}           & Zoom Out \\
  \key{Ctrl} + \key{0}           & Reimposta Zoom (100\%) \\
  \key{Delete} / \key{Backspace} & Elimina elemento selezionato \\
  \key{Esc}                      & Deseleziona / Passa alla modalità Seleziona \\
  \bottomrule
\end{tabular}
\end{center}

\begin{tcolorbox}[notebox, title={Nota}]
  La scorciatoia \key{Ctrl} + \key{S} apre automaticamente la finestra ``Salva con nome''
  se il diagramma non è ancora stato salvato su file.
\end{tcolorbox}

% =============================================================
\section{Formato dei File}
\label{sec:files}

\jder{} salva i diagrammi in formato \textbf{JSON} (JavaScript Object Notation). I file sono
leggibili e modificabili con qualsiasi editor di testo, sebbene la modifica manuale non sia
consigliata.

\begin{center}
\begin{tabular}{@{}lll@{}}
  \toprule
  \textbf{Tipo di diagramma} & \textbf{Estensione} & \textbf{Contenuto principale} \\
  \midrule
  Entità-Relazione    & \texttt{.json} & entità, relazioni, attributi, note \\
  Diagramma Classi    & \texttt{.json} & classi, relazioni, note \\
  Casi d'Uso          & \texttt{.json} & attori, casi d'uso, relazioni, sistemi, note \\
  \bottomrule
\end{tabular}
\end{center}

Tutti e tre i formati condividono la stessa estensione \texttt{.json} ma hanno strutture interne
differenti (non sono intercambiabili). La serializzazione è gestita da \texttt{kotlinx.serialization};
ogni oggetto è identificato da un UUID univoco generato automaticamente.

% =============================================================
\section{Risoluzione dei Problemi}
\label{sec:troubleshooting}

\begin{tcolorbox}[warnbox, title={L'applicazione non si avvia}]
  \begin{itemize}
    \item Verificare che Java 17 (o superiore) sia installato: \texttt{java -{}-version}.
    \item Su Windows, provare ad avviare da terminale per leggere il messaggio di errore.
    \item Assicurarsi che il file \texttt{JDER.jar} non sia bloccato (clic destro
          $\to$ Proprietà $\to$ Sblocca su Windows).
  \end{itemize}
\end{tcolorbox}

\begin{tcolorbox}[warnbox, title={Il file salvato non si apre}]
  \begin{itemize}
    \item Assicurarsi di aprire il file nella scheda corretta (es.\@ un file di diagramma classi
          deve essere aperto nella scheda \uielement{Classi}).
    \item Il file potrebbe essere corrotto. Tentare di aprirlo in un editor di testo
          per verificarne la struttura JSON.
  \end{itemize}
\end{tcolorbox}

\begin{tcolorbox}[warnbox, title={L'esportazione PNG è vuota o incompleta}]
  \begin{itemize}
    \item Verificare che il diagramma contenga elementi visibili sulla lavagna.
    \item Provare a reimpostare lo zoom prima di esportare.
  \end{itemize}
\end{tcolorbox}

% =============================================================
\section{Crediti e Licenza}
\label{sec:crediti}

\jder{} è un progetto open-source rilasciato sotto licenza \textbf{MIT}.

\subsection*{Autori e ringraziamenti}
\begin{itemize}
  \item \textbf{Versione attuale (Kotlin / Compose):} LoryPelli
        (\url{https://github.com/LoryPelli/JDER})
  \item \textbf{Versione precedente Java:} Gianvito Pio
        (\url{https://gianvitopio.wordpress.com/jder})
  \item \textbf{Versione originale:} Alessandro Ballini
        (\url{https://ballini.it/Software/ProgER})
\end{itemize}

\subsection*{Tecnologie utilizzate}
\begin{itemize}
  \item \textbf{Kotlin} --- linguaggio di programmazione principale
  \item \textbf{Jetpack Compose Multiplatform} --- framework UI dichiarativo
  \item \textbf{Material 3} --- sistema di design dell'interfaccia
  \item \textbf{kotlinx.serialization} --- serializzazione JSON
\end{itemize}

\subsection*{Risorse}
\begin{itemize}
  \item Sito web ufficiale: \url{https://jder.it}
  \item Repository GitHub: \url{https://github.com/LoryPelli/JDER}
  \item Segnalazione problemi: \url{https://github.com/LoryPelli/JDER/issues}
  \item Download ultima versione: \url{https://github.com/LoryPelli/JDER/releases/latest}
\end{itemize}

\vfill
\begin{center}
  \textcolor{jderorange}{\rule{8cm}{1pt}}\\[4pt]
  {\small \textit{Guida Utente \jder{} --- \theauthor}}
\end{center}

\end{document}
